\scanpage{14}

\mspage{7)}
qui se composent encore comme précédemment. On a bien \deleted{\unsure{un|sur|donc}} foncteur

\[
 \deleted{\unsure{C\longrightarrow}}\quad\widetilde C\longrightarrow W^{-1}C
\]

% slanted
\margin{Ce foncteur est bijectif sur les objets, surjectif sur les flèches, comme on vérifie aisément (il suffit pour ce dernier point de voir que pour tt $f\in W$, $f:x\longrightarrow y$, $[f]^{-1}$ provient de $\widetilde C$, \unsure{or prenant|on prend} la flèche \unsure{idem|'id'}.)}

\diagram{\begin{tikzcd}[ampersand replacement=\&,column sep=small,row sep=small]
x \arrow[r,"i"] \arrow[d,"\sim"'] \& y \\
\deleted{\unsure{y}} \arrow[ur,"{\alpha}"] \& {}
\end{tikzcd}}

\diagram{\begin{tikzcd}[ampersand replacement=\&,column sep=small,row sep=small]
y \arrow[dr,"{1_{\unsure{x|y}}}"'] \& x \arrow[d,"f","\sim"'] \\
{} \& y
\end{tikzcd}}

\deleted{\unsure{Le premier|Le processus} de ces deux foncteurs transforme bien éléments de $W$ en flèches inversibles. \deleted{On a} Pour $f:x\xrightarrow[\sim]{f}y$ dans $W$, on considère soit $F=[f]$}

\diagram{\begin{tikzcd}[ampersand replacement=\&,column sep=small,row sep=small]
\deleted{y} \arrow[dr,tail,"{\deleted{1_y}}"'] \& \deleted{x} \arrow[d,"{\deleted{f}}","{\deleted{\sim}}"'] \\
{} \& \deleted{y}
\end{tikzcd}}

\deleted{soit $G=[1_y,f]$ on trouve les composés \deleted{On a donc} $GF=[\quad{}^x$ définissons de $\mathord{=}$}

\[
 C\longrightarrow\widetilde C
\]

qui va être un peu plus délicat. Soit $f:x\longrightarrow y$ on choisit une factorisation

\diagram{\begin{tikzcd}[ampersand replacement=\&,column sep=small,row sep=small]
a \arrow[dr,tail,"{f_0}"'] \& y \arrow[dr,dotted,bend left,"i"] \& {} \\
{} \& \bar y \arrow[u,"\alpha","\sim"'] \arrow[r,tail,dotted,"{\sim\,j}"'] \& \bar y' \arrow[ul,dashed,"\beta","\sim"']
\end{tikzcd}}

donc $\bar y$ cof

et : en vertu du weak lifting lemma (et $y$ étant cofibrant) on peut compléter par des pointillés, où $\beta$ a une section $i$. On considère alors $[jf_0,i]$, il faut montrer que ça ne dépend pas des choix faits. Donc on choisit $[f',i]$

\margin{\unsure{Ou s'est ramené aux|On s'est ramené aux} factorisations $f=\deleted{\unsure{\beta}}\alpha f_0$, \inserted{avec} $f_0\in\mathrm{cof}$ et $\deleted{\unsure{\beta}}\alpha\in W$, telles que $\deleted{\unsure{\beta}}\alpha$ ait une section}

\diagram{\begin{tikzcd}[ampersand replacement=\&,column sep=small,row sep=small]
a \arrow[r,dashed,"f"] \arrow[dr,tail,"{f'}"'] \& y \arrow[d,"i","\sim"'] \\
{} \& \unsure{\widetilde y|y'} \arrow[u,bend right,dashed,"\alpha","\sim"']
\end{tikzcd}}

tel qu'il existe \deleted{\unsure{g}} une rétraction $\alpha$ de $\widetilde y$ sur $y$, telle que $\alpha f'=f$

Considérons deux \deleted{\unsure{tels|telles}} choix, il faut voir qu'ils sont dans une $=$ composante connexe de $\underline{\mathrm{Ch}}_0(x,y)$. Mais pour deux factorisations d'une flèche $f:x\longrightarrow y$ en $\alpha f_0$ et $\alpha'f'_0$, $f_0,f'_0\in\mathrm{cof}$, $\alpha,\alpha'\in W$,

\deleted{$x\longrightarrow\unsure{y}$}

On peut les coiffer par une troisième \deleted{\unsure{$(\alpha'',f)$}} $\alpha'',f''_0$ qui \deleted{provient de} \inserted{recolle} les deux comme on voit sur le diagramme

\diagram{\begin{tikzcd}[ampersand replacement=\&,column sep=small,row sep=small]
\deleted{a} \arrow[r,"\deleted{\unsure{f}}"] \arrow[dr] \& \deleted{y} \arrow[d] \\
{} \& \deleted{y''}
\end{tikzcd}}

% archive sheet number: 13
